% Source: https://artofproblemsolving.com/community/c3h2622720_symmetry_in_absolute_value_equation
% Stable identity: aops-topic-2622720-q01
Let $f(x)$ be the left side. Replacing $x$ by $-7-x$ permutes the four terms ($|-7-x+2|=|x+5|$, and so on), so $f(-7-x)=f(x)$. The graph is symmetric about $x=-\dfrac72$, and it is enough to solve for $x\ge-\dfrac72$.
\begin{itemize}
\item For $-\dfrac72\le x\le-3$: $f(x)=-(x+2)-(x+3)+(x+4)+(x+5)=4\ne6$.
\item For $-3\le x\le-2$: $f(x)=-(x+2)+(x+3)+(x+4)+(x+5)=2x+10$, which equals $6$ only at $x=-2$.
\item For $x\ge-2$: $f(x)=4x+14$, which equals $6$ only at $x=-2$.
\end{itemize}
So $x=-2$ is the only solution with $x\ge-\dfrac72$, and by symmetry the other is $-7-(-2)=-5$. Check: $f(-2)=0+1+2+3=6$ and $f(-5)=3+2+1+0=6$. The solutions are $\boxed{x=-5\text{ and }x=-2}$.
